Mã hóa khóa công khai giải quyết một bài toán mà mã hóa đối xứng đơn lẻ xử lý chưa thuận tiện: thiết lập khóa an toàn giữa các bên chưa chia sẻ bí mật từ trước. Trong chương này, hai cấu trúc lớn minh họa cách mục tiêu đó được đạt tới trên những nền tảng toán học khác nhau.

Cấu trúc thứ nhất, \textbf{ElGamal}, dựa trên phép lũy thừa trong một nhóm cyclic và nhận độ an toàn từ độ khó của các bài toán kiểu Diffie--Hellman, đặc biệt là \textbf{giả thuyết Decisional Diffie--Hellman}. Việc sử dụng ngẫu nhiên mới trong mỗi lần mã hóa là yếu tố thiết yếu, bởi chính thành phần ngẫu nhiên này bảo đảm rằng bản mã không để lộ tính bằng nhau giữa các thông điệp và hỗ trợ semantic security.

Cấu trúc thứ hai, \textbf{RSA}, được xây dựng từ một trapdoor permutation trên số học modulo một số hợp thành. Tính đúng đắn của nó suy ra từ định lý Euler, còn an toàn thực tiễn của nó gắn với độ khó của việc phân tích thừa số của mô-đun. RSA cho thấy cách một cấu trúc lý thuyết số thanh nhã có thể tạo ra một nguyên thủy khóa công khai hiệu quả với chiều thuận rõ ràng và chiều đảo ngược được che giấu.

Đồng thời, chương này cũng chỉ ra rằng \textbf{textbook RSA là không an toàn}. Tính tất định, điểm yếu trước thông điệp nhỏ và cấu trúc nhân của RSA thô đều ngăn cản nó trở thành một sơ đồ mã hóa an toàn nếu đứng riêng lẻ. Vì vậy, một trapdoor permutation không đồng nghĩa với một hệ mã khóa công khai an toàn.

Để đạt được an toàn trong thực tế, RSA phải được kết hợp với \textbf{padding}. Padding đưa ngẫu nhiên và cấu trúc tiền xử lý vào trước khi thực hiện phép lũy thừa, nhờ đó loại bỏ những điểm yếu rõ ràng nhất của textbook RSA. Các cách mã hóa cũ như PKCS\#1 v1.5 đã cải thiện đáng kể tình hình, nhưng lịch sử cho thấy chúng vẫn có thể thất bại dưới các điều kiện tấn công bản mã được chọn. Những thiết kế hiện đại hơn như OAEP được đưa ra để cung cấp mức bảo vệ mạnh hơn.

\textbf{Tấn công Bleichenbacher} là ví dụ thực tế rõ ràng nhất cho nguyên lý này. Nó cho thấy một hệ mật mã có thể bị đánh bại không phải bằng cách phá vỡ lõi toán học của nó, mà bằng cách khai thác thông tin rò rỉ trong quá trình cài đặt. Theo nghĩa đó, bài học quan trọng nhất của chương này vượt ra ngoài riêng ElGamal hay RSA: \textbf{an toàn mật mã phụ thuộc đồng thời vào lý thuyết và cài đặt}. Các giả thuyết độ khó mạnh là cần thiết, nhưng chúng không đủ nếu toàn bộ sơ đồ không được thiết kế và triển khai với mức độ cẩn trọng tương xứng.
